\documentclass[11pt]{article}
\newcommand{\DocumentKind}{Manuscript}
\usepackage[a4paper,top=25mm,bottom=25mm,left=28mm,right=28mm,headheight=15pt,footskip=11mm]{geometry}
\usepackage{fontspec}
\setmainfont{Latin Modern Roman}
\setsansfont{Latin Modern Sans}
\setmonofont{Latin Modern Mono}
\usepackage{amsmath,amssymb,amsthm,mathtools,bm}
\usepackage{booktabs,longtable,tabularx,array}
\usepackage{enumitem,graphicx,listings,fancyhdr,needspace}
\usepackage{url}
\usepackage[hidelinks,unicode]{hyperref}
\usepackage{xr}
\setlength{\parindent}{1.2em}
\setlength{\parskip}{0pt}
\setlength{\emergencystretch}{3em}
\renewcommand{\baselinestretch}{1.08}
\raggedbottom
\setlist[enumerate]{leftmargin=2em,itemsep=0.15em,topsep=0.35em}
\setlist[itemize]{leftmargin=1.8em,itemsep=0.15em,topsep=0.35em}
\newtheorem{theorem}{Theorem}[section]
\newtheorem{lemma}[theorem]{Lemma}
\newtheorem{proposition}[theorem]{Proposition}
\newtheorem{corollary}[theorem]{Corollary}
\theoremstyle{definition}
\newtheorem{definition}[theorem]{Definition}
\newtheorem{example}[theorem]{Example}
\theoremstyle{remark}
\newtheorem{remark}[theorem]{Remark}
\numberwithin{equation}{section}
\newcommand{\Nout}{N^+}
\newcommand{\Nin}{N^-}
\newcommand{\Ntwo}{N^{++}}
\newcommand{\Nback}{N^{--}}
\newcommand{\dout}{d^+}
\newcommand{\din}{d^-}
\newcommand{\dtwo}{d^{++}}
\newcommand{\dback}{d^{--}}
\newcommand{\bd}{\partial}
\newcommand{\ind}{\mathbf 1}
\DeclareMathOperator{\Reach}{Reach}
\DeclareMathOperator{\med}{med}
\pagestyle{fancy}
\fancyhf{}
\fancyhead[L]{\small\itshape Incoming cores and extremal constraints}
\fancyhead[R]{\small\DocumentKind}
\fancyfoot[C]{\small\thepage}
\renewcommand{\headrulewidth}{0.3pt}
\renewcommand{\footrulewidth}{0pt}
\setcounter{tocdepth}{1}
\lstdefinestyle{journalcode}{basicstyle=\footnotesize\ttfamily,keywordstyle=\rmfamily\bfseries,
commentstyle=\itshape,showstringspaces=false,frame=none,numbers=left,
numberstyle=\scriptsize,numbersep=7pt,xleftmargin=1.5em,breaklines=true,
columns=fullflexible,keepspaces=true,tabsize=2,captionpos=b}
\lstset{style=journalcode}
\lstdefinelanguage{JavaScript}{morekeywords={const,let,var,function,return,for,of,in,if,else,new,throw,true,false,null,require},sensitive=true,morecomment=[l]{//},morecomment=[s]{/*}{*/},morestring=[b]",morestring=[b]'}
\makeatletter
\renewcommand*\l@section{\@dottedtocline{1}{0em}{3.4em}}
\makeatother
\def\UrlBreaks{\do\/\do-\do_\do\%\do.}
\hypersetup{pdfauthor={Guanze Yu},pdfsubject={Research preprint; complete conjecture not resolved}}

\externaldocument[S-]{../supplement/main}[Seymour_Supplement_2026-10-07.pdf]
\title{Five Maximal Incoming Classes Suffice:\par
Exact Certificates and Structural Constraints for\par
Seymour's Second-Neighbourhood Conjecture}
\author{Guanze Yu}
\date{7 October 2026}
\hypersetup{pdftitle={Five Maximal Incoming Classes Suffice: Exact Certificates and Structural Constraints for Seymour's Second-Neighbourhood Conjecture},pdfauthor={Guanze Yu}}
\begin{document}
\maketitle
\begin{abstract}
We study Seymour's second-neighbourhood conjecture through incoming-neighbourhood
patterns, counterexample-preserving row maps, and protection-preserving tournament
completions. A finite system of homogeneous pattern inequalities provides an
exact-second-neighbourhood certificate in the original graph. We prove that
universal feasibility of these systems is equivalent to the conjecture, and
verify feasibility for every oriented core on at most five vertices using
nonnegative integer certificates. Consequently, every finite nonempty oriented graph
with at most five inclusion-maximal incoming-neighbourhood classes has a whole
maximal class of Seymour vertices, with no bound on its total order.
We also prove a simultaneous outgoing-row copying lemma, obtain an independent-twin
normal form for a lexicographically extremal counterexample, and locate a global
core of gap-one maximal incoming classes. Iterated protection closure preserves
the original first and exact second neighbourhoods at its partial median feeds.
The resulting tournament feeds have a strict original far-head surplus; at zero
completed gap, a safe weak closure forces an original cycle among the remaining
missing partners. In the arc-set-minimal setting this gives residual at least four
at every such zero-gap feed. Together with the supplementary residual exclusions,
we obtain the necessary bound $2M-n\ge7$, strengthened to eight at even order,
where $n$ is the order and $M$ is the number of missing unordered vertex pairs
in the arc-set-minimal counterexample under consideration.
All finite certificates and independent integer verifiers are supplied. These
results do not prove or refute the conjecture in full generality.
\end{abstract}
\noindent\textbf{Keywords:} oriented graph; second neighbourhood; tournament;
median order; incoming-neighbourhood inclusion; computer-assisted proof.
\medskip
\noindent\textbf{MSC 2020:} 05C20, 05C22, 05C85.
\section{Introduction}
For a vertex $v$ of a finite oriented graph $D$, let $N_D^+(v)$ denote its
outneighbourhood and let $N_D^{++}(v)$ consist of vertices whose directed
distance from $v$ is exactly two. Seymour's conjecture asks whether every
nonempty oriented graph has a vertex satisfying
\begin{equation}\label{eq:paper-snc}
 |N_D^{++}(v)|\ge |N_D^+(v)|.
\end{equation}
Such a vertex is called a Seymour vertex. The exclusion of first neighbours
from the second set is essential in every deletion, completion and counting
argument below.

The tournament case is known, but a general completion may create second
neighbours that were absent in the original graph. We develop two ways of
making that discrepancy explicit. The first uses a small induced incoming
core and an exact linear certificate. The second completes only carefully
controlled directions before a global median-order optimization. A separate
copying operation imposes global structural restrictions on a suitably chosen
counterexample. The three constructions have different preservation properties;
their hypotheses are kept distinct throughout the paper.

\subsection{Principal results}
Partition vertices by equality of their full incoming neighbourhoods. A class
is maximal when its common incoming set is maximal under inclusion among the
sets $N_D^-(v)$. Maximality refers to inclusion, not to the numerical indegree.
Our first unconditional result is the following computer-assisted theorem.

\begin{theorem}[Five maximal incoming classes]\label{thm:intro-five}
Every finite nonempty oriented graph with at most five maximal incoming-neighbourhood
classes has a maximal class all of whose vertices are Seymour vertices.
The total order and the sizes of the classes are unrestricted.
\end{theorem}

The theorem follows from a transfer inequality valid for arbitrary induced
cores. Its finite component comprises $59\,809$ labelled cores of orders one
through five and $462\,404$ nonempty pattern inequalities. Every witness is a
nonnegative integer vector, with entries at most three. The tables are checked
independently of the program that found them. The full proof and the precise
finite verification contract appear in Theorem~\ref{thm:core-five}.

The same certificate system has a converse: an infeasible core system yields,
by a finite labelled construction, a genuine all-deficient oriented graph
(Theorem~\ref{thm:pattern-duality}). Thus universal feasibility is equivalent
to the full conjecture. A certificate also lifts through repeated maximal
incoming-core reductions, reducing the remaining feasibility question to
cores with pairwise distinct, incomparable incoming sets.

Our second group of results concerns simultaneous copying. If an original
vertex map $\rho$ satisfies $N_D^-(v)\subseteq N_D^-(\rho(v))$ for every $v$,
replacing each outgoing row by the original row of $\rho(v)$ creates an
oriented graph whose gap at $v$ is at least the original gap at $\rho(v)$.
This implies a global covering property in a smallest-order counterexample:
every vertex receives an entire maximal incoming class whose members all have
gap one. A further extremal selection makes equal incoming neighbourhoods into
independent true-twin modules. These statements use global extremality rather
than merely inclusion-minimal deletion assumptions.

For the completion results write $g_D(v)=d_D^+(v)-d_D^{++}(v)$ and let $t_D(v)$
count the original far vertices not protected by incoming-neighbourhood nesting.
The closure constructed in Section~\ref{sec:paper-closure} admits an optimal
reference completion whose every original-positive-gap median feed satisfies
\begin{equation}\label{eq:intro-surplus}
 t_D(f)\ge g_D(f)+1.
\end{equation}
Under the safe weak extension, zero completed gap requires at least three
remaining missing partners, an original directed cycle among them, and
$t_D(f)\ge g_D(f)+3$. For a strongly connected all-deficient graph minimal
under every nonempty arc-set deletion, the same feed satisfies
$\eta_D(f)\ge4$, where
$\eta_D(v)=\mu_D(v)-1+d_D^{--}(v)-d_D^{++}(v)$.

The supplementary case reductions, together with an unbounded one-exception
theorem for actual bad-direction failure sources, yield
\begin{equation}\label{eq:intro-residual}
 2M-n\ge7,\qquad
 M\ge\left\lceil n/2\right\rceil+3+\ind_{\{n\text{ even}\}}.
\end{equation}
The missing-pair bound is a necessary condition for the stated arc-set-minimal
class. It is not a missing-edge density theorem for arbitrary original graphs:
deleting arcs changes $M$. Proofs of the residual allocations and the specialized
tools for $n=2\delta^++3$ are collected in the companion supplement.

\subsection{Relation to previous work}
\label{subsec:related-work}

The tournament case of Seymour's second neighborhood conjecture was proved
by Fisher~\cite{Fisher1996} using a probability distribution obtained from
a skew-symmetric game. Havet and Thomass\'e~\cite{HavetThomasse2000}
subsequently gave a proof using median orders, including a stronger bound
involving the good vertices of an order. Both methods enter the present
work. The incoming-core argument uses the existence of a losing density
on a tournament; the completion arguments use the strong feed bound and
the equality case of sedimentation. These are established inputs. Our
arguments specify the additional conditions needed to transfer information
from a tournament completion to the exact second neighborhoods of the
original graph. In particular, the existence of a Seymour vertex in a
completion alone does not provide such a transfer.

Fidler and Yuster~\cite{FidlerYuster2007} proved the vertex-weighted
tournament statement and established further cases obtained by deleting
prescribed edge sets from tournaments. Ghazal's work on generalized stars
and other missing graphs~\cite{Ghazal2012,Ghazal2013} develops the convenient
orientation and good-missing-edge approach underlying several of our
completion definitions. Dara, Francis, Jacob and
Narayanan~\cite{DaraEtAl2022} extended the matching case to a missing graph
whose edges can be partitioned into a matching and a star, and also
considered partitions into an independent set and a $2$-degenerate graph.
Our protection-preserving completions require additional compatibility
conditions beyond convenience of individual added arcs. The proof of each
transfer inequality therefore accounts for paths using added arcs and for
original second neighbors that become first neighbors. None of these
inequalities is inferred simply by applying a known tournament theorem to
an arbitrary completion.

Minimal-counterexample methods provide another starting point.
Brantner, Brockman, Kay and Snively~\cite{BrantnerEtAl2009} established
strong connectivity, gaps in $\{1,2\}$, and restrictions involving
gap-one predecessors and directed cycles for their minimal counterexamples.
Their extremal selection first minimizes the number of arcs over all
counterexamples and then minimizes the order. The selections used here
must be distinguished from that convention: fixed-order minimality under
deletion of a nonempty arc set, and minimum order followed by arc count
and further finite objectives, serve different arguments. We give the
required deletion and copying arguments for the stated selection in each
case. In particular, the predecessor-containment row operation and the
resulting normal form are not treated as consequences of the earlier
extremal definition.

Seacrest's subset formulation~\cite{SeacrestSubsets2019} supplies a closely
related boundary perspective. Its edge-minimal subset inequality motivates
the boundary budgets used in the residual analysis. Because the sets and
minimality assumptions have to agree exactly, the external-boundary form
needed here is proved explicitly. This paper is distinct from Seacrest's
arc-weighted extension~\cite{Seacrest2015}; an arc-weighted result is not
substituted for the unweighted feed or boundary statements. The residual
bound $R\geq 7$ in this manuscript concerns the specified minimal
all-deficient graphs. It must not be read as an assertion about the
missing-edge count of every counterexample before a minimizing deletion.

There are further reasons to keep extremal assumptions visible.
Espuny D\'iaz, Gir\~ao, Granet and
Kronenberg~\cite{EspunyDiazEtAl2025} prove that a vertex-minimal
counterexample has minimum outdegree greater than the square root of its
order, while also constructing arbitrarily large strongly connected
counterexamples with bounded minimum outdegree, conditional on the
existence of any counterexample. Zelenskiy, Darmosiuk and
Nalivayko~\cite{ZelenskiyEtAl2021} obtain conditional counterexample
constructions with widely differing densities and diameters. Thus
counterexample amplification, and cyclic substitution in particular,
already have a literature. The calculations here involving such
constructions are used to delimit particular completion or weighted
transfer claims, rather than to assert a new general amplification
principle.

Theorem~\ref{thm:core-five} has a different parameter from these density
and degree results: the number of distinct inclusion-maximal incoming
neighborhoods. There is no restriction on the order of the original graph
or on the size of an incoming-neighborhood class. The certificate on a
chosen core is checked against every permitted subset of an incoming set,
so that it controls all original target patterns simultaneously. An
ordinary verification of Seymour's conjecture for small graphs does not
automatically supply these stronger certificates. The finite tables are
therefore accompanied by explicit integer checks and a proof of their
transfer to arbitrary-order graphs. The labelled construction from an
infeasible pattern system establishes a converse at the universal level.
Finite linear-programming duality is standard; the graph realization and
its exact-distance calculation are the substantive steps that must be
verified for this particular reformulation. Universal feasibility remains
unproved.

Recent special cases provide useful comparisons. Wang and
Lu~\cite{WangLu2026} prove the conjecture for
$\mathcal D_{0,3}\cup\mathcal D_{1,1}$, where $\mathcal D_{s,t}$ consists
of oriented graphs admitting a vertex partition whose induced underlying
graphs are respectively $s$- and $t$-degenerate. Bai, Li and
Park~\cite{BaiLiPark2026} study the stronger requirement of a complete
directed matching from the first to the exact second neighborhood, proving
it when the minimum outdegree is at most five or the graph is
$5$-anti-transitive. Kaneko and Locke~\cite{KanekoLocke2001} established
ordinary Seymour vertices when the minimum outdegree is at most six.
The preprint of Sadhukhan, Sandeep and Sen~\cite{SadhukhanEtAl2026}
treats minimum outdegree seven using computer-assisted local reductions.
Brukhman's preprint~\cite{Brukhman2026} treats the dense range
$n\leq 2\delta^++2$. The incoming-core theorem uses none of these
recent preprints as a premise.

The five-class hypothesis does not imply the hypotheses of those
particular recent results. For example, replace each vertex of the regular
five-vertex tournament, with arcs $i\to i+1,i+2$ modulo five, by an
independent set of size $m\geq4$. The resulting graph has five maximal
incoming classes, order $5m$, and minimum outdegree $2m$. Its underlying
complete five-partite graph belongs to neither
$\mathcal D_{0,3}$ nor $\mathcal D_{1,1}$, and a path through parts
$0,1,2,3,4,1$ excludes $5$-anti-transitivity. It also lies outside the
stated dense range. This example itself is already covered by the weighted
tournament theorem. It demonstrates noncontainment in specific hypotheses,
not novelty relative to every known weighted or substitution argument.

Halkiewicz's split-twin operation~\cite{Halkiewicz2026} duplicates incoming
neighbors and partitions outgoing neighbors while increasing the order;
its stated preservation theorem has a condition involving an existing
Seymour vertex. This differs from the fixed-order row copying used here.
Nevertheless, this distinction and the preceding comparisons do not settle
priority for equivalent formulations under other terminology. The present
paper makes no claim to an exhaustive classification of earlier reductions.
The literature comparison is current to 7 October 2026. Public full-proof
claims, including~\cite{GloverClaim2026}, are not assumed correct or used
as premises. The conclusions proved here remain partial results and
reformulations of the conjecture.


\subsection{Organization and verification}
The incoming-core and cloning arguments are independent of the residual
machinery. After establishing them, we give the closure construction and the
zero-gap cycle argument, then state the minimal-counterexample consequences
with their supplementary proof locations. The final sections give reproducibility
details, exact controls against stronger transfer claims, and the unresolved
global steps. The supplement is a self-contained technical repository of the
longer foundational and case arguments. It also records which earlier finite
classifications are subsumed by the stronger results.

\section{Notation and extremal hypotheses}\label{sec:paper-notation}
All graphs are finite and nonempty. An oriented graph has no loops and no
pair of oppositely directed arcs. Set
\[
 N_D^{++}(v)=\left(\bigcup_{a\in N_D^+(v)}N_D^+(a)\right)
       \setminus\bigl(N_D^+(v)\cup\{v\}\bigr),
 \qquad g_D(v)=d_D^+(v)-d_D^{++}(v).
\]
Reverse exact second neighbourhoods are denoted by $N_D^{--}(v)$.
We call $D$ all-deficient when $g_D(v)>0$ for every vertex. A missing
partner of $v$ is a different vertex nonadjacent to it; its set and size
are $\mathcal M_D(v)$ and $\mu_D(v)$. The number of missing unordered
pairs is $M(D)$, so $\sum_v\mu_D(v)=2M(D)$.

The far set and the reverse far set are
\[
 U_D(v)=V(D)\setminus\bigl(\{v\}\cup N_D^+(v)\cup N_D^{++}(v)\bigr),
 \qquad P_D(v)=\{x\ne v:v\in U_D(x)\}.
\]
For $S\subseteq V(D)$ we use external boundaries
\[
 \Gamma_D(S)=\bigcup_{x\in S}N_D^+(x),\qquad
 \partial_D S=\Gamma_D(S)\setminus S,\qquad
 \partial_D^2S=\Gamma_D(\partial_D S)\setminus(S\cup\partial_D S).
\]
The word external distinguishes these sets from conventions that allow a
set neighbourhood to meet the starting set.

We will use three different extremal choices. An all-deficient graph is
\emph{arc-set-minimal} when deleting every nonempty set of arcs on its fixed
vertex set destroys all-deficiency. The abbreviation MIN means, in addition,
that it is strongly connected. A \emph{minimum-order} counterexample has the
smallest order among all counterexamples, including graphs obtained by
rewiring. The lexicographic normal form in Section~\ref{sec:paper-predecessor-cloning}
adds further global objectives. These conditions are not interchangeable.
If a counterexample exists, a MIN graph exists by taking an arc-minimal
spanning subgraph and then a sink strong component. A minimum-order
counterexample is automatically strongly connected.

\begin{definition}[Protection]\label{def:paper-protection}
Write $a\triangleleft_D b$ if $a\to b$ and
$N_D^-(a)\subseteq N_D^-(b)\setminus\{a\}$.
Then $a$ is a protected far head of $b$. This is a strict partial order.
Put
\[
 \mathcal F_D(b)=U_D(b)\setminus\{a:a\triangleleft_D b\},
 \qquad t_D(b)=|\mathcal F_D(b)|.
\]
\end{definition}
Indeed, an alleged path $b\to x\to a$ would force $x\to b$, a digon.
Transitivity follows because $a\to b\to c$ and predecessor containment
give $a\to c$ and the required further containment. A whole vertex, meaning
$\mu_D(v)=0$, has every far head protected.

A direction $a\to b$ on a missing pair is \emph{convenient} if every
original predecessor of $a$ already reaches $b$ in one or two steps.
A missing pair is good if at least one direction is convenient, and bad
otherwise. All such predicates carry an explicit graph index whenever the
graph has changed. A \emph{global median order} maximizes the total number
of forward arcs among all vertex orders. Its last vertex is its feed.
Local feedback inequalities alone do not justify the sedimentation step
used later.

\begin{remark}[Two kinds of unit]
Gap one means $g_D(v)=1$. Residual one means $\eta_D(v)=1$.
Actual failure-source status refers to a witnessed nonconvenient direction
of an original bad pair. None of these three notions is identified with
another without an explicit lemma.
\end{remark}

\section{Incoming-neighborhood pattern certificates}
\label{sec:paper-incoming-certificates}

We first give a certificate that depends on a small induced subgraph but
controls second neighborhoods in the original graph. Throughout this section,
$D$ is a finite nonempty oriented graph. Put
\[
 g_D(v)=|N_D^+(v)|-|N_D^{++}(v)|,
\]
where $N_D^{++}(v)$ consists of the vertices at directed distance exactly
two from $v$. A vertex with $g_D(v)\leq0$ is called a Seymour vertex.
No extremal or minimal-counterexample assumption is made in this section.

\subsection{Patterns and the transfer theorem}

Let $H$ be a finite nonempty oriented graph. Define its family of
\emph{permitted incoming patterns} by
\begin{equation}
 \mathcal P(H)=\{P\subseteq V(H):
           P\subseteq N_H^-(j)\text{ for some }j\in V(H)\}.
 \label{eq:paper-pattern-family}
\end{equation}
In particular, the empty set belongs to $\mathcal P(H)$. For an arbitrary
$P\subseteq V(H)$, write
\[
 R_H(P)=\{i\in V(H)\setminus P:
                      N_H^+(i)\cap P\ne\varnothing\}.
\]
Thus $R_H(P)$ consists of the vertices outside $P$ that send an arc into
$P$. If $\lambda$ is a nonnegative vector indexed by $V(H)$, use the
notation $\lambda(S)=\sum_{i\in S}\lambda_i$.

\begin{definition}[Pattern certificate]
\label{def:paper-pattern-certificate}
A pattern certificate for $H$ is a nonzero vector $\lambda\geq0$ such that
\begin{equation}
 \lambda(R_H(P))\geq\lambda(P)
             \qquad\text{for every }P\in\mathcal P(H).
 \label{eq:paper-pattern-inequalities}
\end{equation}
The inequalities are homogeneous, so a certificate may always be
normalized by $\lambda(V(H))=1$.
\end{definition}

For a nonempty set $C\subseteq V(D)$, let $H=D[C]$ and associate to every
original vertex $v$ the pattern
\[
 \Pi_C(v)=N_D^-(v)\cap C.
\]
We say that $C$ is \emph{incoming-covered} if
\begin{equation}
 \Pi_C(v)\in\mathcal P(H)\qquad(v\in V(D)).
 \label{eq:paper-incoming-covered}
\end{equation}
The witness for this condition may depend on $v$; the condition does not
require any particular orientation between $v$ and that witness.

For clarity, let $\mathbf N_D$ denote the matrix with rows and columns
indexed by $V(D)$ and entries
\[
 (\mathbf N_D)_{uv}=
 \begin{cases}
 -1,&v\in N_D^+(u),\\
  1,&v\in N_D^{++}(u),\\
  0,&\text{otherwise}.
 \end{cases}
\]
The row sum at $u$ is $-g_D(u)$.

\begin{theorem}[Transfer from an incoming-covered core]
\label{thm:pattern-transfer}
Let $C\subseteq V(D)$ be nonempty and incoming-covered, and put $H=D[C]$.
If $\lambda$ is a pattern certificate for $H$, then its extension
$\widehat\lambda$ by zero to $V(D)$ satisfies
\[
 \mathbf N_D^{\mathsf T}\widehat\lambda\geq0.
\]
Consequently some $c\in C$ with $\lambda_c>0$ is a Seymour vertex of $D$.
\end{theorem}
\begin{proof}
Fix an original target $v$ and put $P=\Pi_C(v)$. The first predecessors of
$v$ in $C$ are precisely the vertices of $P$. If $i\in R_H(P)$, choose
$j\in N_H^+(i)\cap P$. Both arcs of the path $i\to j\to v$ are
original arcs of $D$, and $i\notin P$ says that $i\to v$ is not an arc.
Also $i\ne v$: otherwise $v\to j$ and $j\to v$ would form a digon.
It follows that $v\in N_D^{++}(i)$.

The total $\lambda$-mass of the core vertices that have $v$ as an exact
second neighbor is therefore at least $\lambda(R_H(P))$. Since $P$ is
permitted, this is at least $\lambda(P)$, the total mass of its core
first predecessors. This proves the asserted matrix inequality at the
column indexed by $v$.

Summing over the columns gives
\[
 0\leq\mathbf1^{\mathsf T}\mathbf N_D^{\mathsf T}\widehat\lambda
       =-\sum_{c\in C}\lambda_c g_D(c).
\]
Since $\lambda$ is nonzero and nonnegative, not every vertex of its
support can have positive gap.
\end{proof}

The certificate weights the possible sources in $C$. The neighborhoods
and their cardinalities in the conclusion remain those of the unweighted
original graph. In particular, vertices outside $C$ are retained as
targets in every column of the transfer inequality.

\subsection{Tournament cores}

Tournament cores admit certificates of arbitrary order. We use the
losing-density formulation underlying Fisher's proof for tournaments
\cite{Fisher1996}. The short argument below also specifies the sign
convention for the skew matrix.

\begin{lemma}[A skew-game density]
\label{lem:paper-skew-density}
Let $B$ be a real skew-symmetric matrix. There is a nonnegative vector
$\lambda$ with $\mathbf1^{\mathsf T}\lambda=1$ and $B\lambda\geq0$.
\end{lemma}
\begin{proof}
Let $\Delta$ be the probability simplex of the appropriate dimension.
Finite minimax gives
\[
 \max_{\lambda\in\Delta}\min_{p\in\Delta}p^{\mathsf T}B\lambda
 =\min_{p\in\Delta}\max_{\lambda\in\Delta}p^{\mathsf T}B\lambda.
\]
The left-hand side is at most zero, by choosing $p=\lambda$ in the
inner minimum. The right-hand side is at least zero, by choosing
$\lambda=p$ in the inner maximum. Both values are therefore zero.
A maximizing $\lambda$ satisfies $p^{\mathsf T}B\lambda\geq0$ for
every $p\in\Delta$, and in particular for every coordinate vector.
\end{proof}

If $A_H$ is the adjacency matrix of a tournament $H$, put
$B_H=A_H-A_H^{\mathsf T}$. A vector supplied by
Lemma~\ref{lem:paper-skew-density} is called a losing density of $H$.
At each vertex, its outgoing density is at least its incoming density.

\begin{lemma}[Expansion of covered patterns in tournaments]
\label{lem:paper-tournament-pattern-expansion}
Every losing density of a nonempty tournament $H$ is a pattern certificate
for $H$.
\end{lemma}
\begin{proof}
Fix $P\subseteq N_H^-(j)$ and set
\[
 R=R_H(P),\qquad F=V(H)\setminus(P\cup R).
\]
The witness $j$ belongs to $F$, so $F$ is nonempty. Tournament
completeness implies that every arc between $P$ and $F$ is directed from
$P$ to $F$: a vertex of $F$ has no arc into $P$ by its definition.

First suppose that $\lambda(F)>0$. Sum the nonnegative quantities
$\lambda_i(B_H\lambda)_i$ over $i\in F$. The terms internal to $F$
cancel by skew-symmetry. The contribution from pairs with the other
endpoint in $P$ is $-\lambda(F)\lambda(P)$. The contribution from pairs
with the other endpoint in $R$ is at most $\lambda(F)\lambda(R)$.
Consequently
\[
 0\leq\sum_{i\in F}\lambda_i(B_H\lambda)_i
 \leq\lambda(F)\bigl(\lambda(R)-\lambda(P)\bigr),
\]
and division by $\lambda(F)$ proves the required inequality.

If $\lambda(F)=0$, apply $(B_H\lambda)_j\geq0$ at any $j\in F$.
The incoming mass at $j$ is at least $\lambda(P)$, whereas its outgoing
mass is at most $\lambda(R)$. This again gives
$\lambda(R)\geq\lambda(P)$. The argument also covers $P=\varnothing$.
\end{proof}

\begin{corollary}[Incoming-covered tournament cores]
\label{cor:paper-tournament-core}
If a nonempty incoming-covered set $C$ induces a tournament in $D$, then
$C$ contains a Seymour vertex of $D$. More precisely, the zero extension
of every losing density of $D[C]$ satisfies the matrix inequality in
Theorem~\ref{thm:pattern-transfer}.
\end{corollary}
\begin{proof}
Combine Lemmas~\ref{lem:paper-skew-density}
and~\ref{lem:paper-tournament-pattern-expansion} with
Theorem~\ref{thm:pattern-transfer}.
\end{proof}

\subsection{Maximal incoming classes and exact small-core certificates}

Define an equivalence relation on $V(D)$ by
\[
 u\sim_D v\quad\Longleftrightarrow\quad N_D^-(u)=N_D^-(v).
\]
The equivalence classes are partially ordered by containment of their
common incoming sets. We call them \emph{incoming classes}, and
maximality will always refer to this inclusion order. Distinct members
of one class cannot be adjacent, since an internal arc would put its tail
in its own incoming set by equality of the two incoming neighborhoods.
An incoming class need not have a common outgoing neighborhood.

\begin{lemma}[Representative cores are incoming-covered]
\label{lem:paper-maximal-representatives-covered}
Choose one representative from each maximal incoming class of $D$, and
let $C$ be the set of chosen vertices. Then $C$ is incoming-covered.
\end{lemma}
\begin{proof}
For every $v\in V(D)$, finiteness of the inclusion poset gives a maximal
class with representative $c$ such that
$N_D^-(v)\subseteq N_D^-(c)$. Intersecting with $C$ gives
\[
 \Pi_C(v)\subseteq N_D^-(c)\cap C=N_{D[C]}^-(c).
\]
This is exactly the required condition.
\end{proof}

\begin{lemma}[Exact certificates for cores of order at most five]
\label{lem:paper-exact-five-certificates}
Every nonempty oriented graph of order at most five has a pattern
certificate. A certificate can be chosen with integer entries in
$\{0,1,3\}$.
\end{lemma}
\begin{proof}[Computer-assisted proof]
The certificate data consist of one nonzero nonnegative integer vector
for every labelled oriented graph of orders one through five. The
encoding lists the unordered pairs
\[
 (0,1),(0,2),\ldots,(0,n-1),(1,2),\ldots,(n-2,n-1)
\]
in lexicographic order. The corresponding ternary digit is zero for a
missing pair, one for the arc from the smaller endpoint to the larger
endpoint, and two for the reverse arc. The first pair corresponds to
the least significant digit. Thus the codes
$0,\ldots,3^{\binom n2}-1$ cover every labelled oriented graph exactly
once.

For every code, the verifier reconstructs its adjacency matrix and checks
that the supplied vector has length $n$, nonnegative integer entries, and
positive sum. It then enumerates every subset $P\subseteq\{0,\ldots,n-1\}$.
A subset is retained precisely when some $j$ satisfies $i\to j$ for
every $i\in P$. For each retained subset, the verifier forms
$R_H(P)$ directly from the reconstructed adjacency matrix and checks
\eqref{eq:paper-pattern-inequalities} with integer arithmetic.

The numbers of graph codes for $n=1,2,3,4,5$ are, respectively,
\[
 1,\quad3,\quad27,\quad729,\quad59049.
\]
All $59809$ certificates pass all permitted-pattern inequalities. There
are $522213$ such checks, of which $462404$ use a nonempty pattern.
The largest integer entry is $3$, and the largest sum of entries is $9$.
Among the order-five certificates, $59009$ are uniform on their supports
and the remaining $40$ are nonuniform. This exhausts the finite assertion.

The complete tables and verification sources accompany the manuscript.
In particular, \texttt{verify\_core\_pattern\_certificates\_independent.js}
reconstructs Boolean adjacency matrices and checks the inequalities
without calling an optimization solver or using a generator's graph
utilities. The integer tables, rather than a floating-point optimization
result, are the certificates used in this proof.
\end{proof}

\begin{theorem}[An entire Seymour maximal class]
\label{thm:core-five}
Every finite nonempty oriented graph with at most five maximal incoming
classes has a maximal incoming class all of whose vertices are Seymour
vertices. There is no restriction on the total number of vertices or on
the sizes of the classes.
\end{theorem}
\begin{proof}
Suppose that no maximal incoming class consists entirely of Seymour
vertices. Choose from each maximal class a vertex with positive gap,
and let $C$ be the set of chosen representatives. Then $1\leq|C|\leq5$,
and Lemma~\ref{lem:paper-maximal-representatives-covered} shows that
$C$ is incoming-covered. Lemma~\ref{lem:paper-exact-five-certificates}
supplies a pattern certificate for $D[C]$.
Theorem~\ref{thm:pattern-transfer} consequently gives a Seymour vertex
in $C$, contrary to the selection of all its members.
\end{proof}

\begin{corollary}
\label{cor:paper-counterexample-six-incoming-classes}
Every counterexample to Seymour's second neighborhood conjecture has at
least six maximal incoming classes.
\end{corollary}
\begin{proof}
An all-deficient graph has no Seymour vertex, so
Theorem~\ref{thm:core-five} applies contrapositively.
\end{proof}

The role of the computation is confined to the finitely many cores in
Lemma~\ref{lem:paper-exact-five-certificates}. The passage to arbitrary
orders and arbitrary class sizes in Theorem~\ref{thm:core-five} uses the
original-graph transfer theorem. It is not an enumeration of all graphs
to which that theorem applies.

\subsection{Strict duality and realization of an infeasible core}

The pattern system also gives an equivalent formulation of the full
conjecture. Its reverse implication requires a realization argument:
infeasibility of a system on a core must yield an actual oriented graph
with positive gap at every vertex.

\begin{theorem}[Pattern feasibility and the second neighborhood conjecture]
\label{thm:pattern-duality}
The following statements are equivalent:
\begin{enumerate}
 \item Every finite nonempty oriented graph has a Seymour vertex.
 \item Every finite nonempty oriented graph $H$ has a pattern certificate.
\end{enumerate}
More explicitly, if the normalized system
\[
 \lambda\geq0,\qquad \sum_i\lambda_i=1,\qquad
 \lambda(R_H(P))\geq\lambda(P)\quad(P\in\mathcal P(H))
\]
is infeasible for some $H$, a finite all-deficient oriented graph can be
constructed from $H$ and a rational strict dual solution.
\end{theorem}
\begin{proof}
First assume statement~(2). For an arbitrary $D$, take the core to be
$D$ itself. Its pattern at a target $v$ is the full set $N_D^-(v)$,
which is permitted. Theorem~\ref{thm:pattern-transfer} yields a Seymour
vertex, proving statement~(1).

For the converse, suppose that the normalized system for a nonempty
core $H$ is infeasible. For each permitted pattern, define the column
vector
\[
 a_P(i)=\mathbf1_{i\in R_H(P)}-\mathbf1_{i\in P}.
\]
On the probability simplex $\Delta$ indexed by $V(H)$, the continuous
function $\min_{P\in\mathcal P(H)}\lambda^{\mathsf T}a_P$ is strictly
negative everywhere. Compactness therefore gives
\begin{equation}
 \max_{\lambda\in\Delta}
       \min_{P\in\mathcal P(H)}\lambda^{\mathsf T}a_P<0.
 \label{eq:paper-strict-dual-value}
\end{equation}
By finite minimax, there are nonnegative real pattern weights whose
weighted sum of the $a_P$ has every coordinate strictly negative.
They may be chosen rational: keep their zero coordinates equal to zero
and approximate the positive coordinates sufficiently closely by
positive rationals, using strictness of all the finitely many
inequalities. Clearing denominators gives nonnegative integers $w_P$
such that
\begin{equation}
 \gamma_i=\sum_{P\in\mathcal P(H)}w_P
       \bigl(\mathbf1_{i\in P}-\mathbf1_{i\in R_H(P)}\bigr)>0
       \qquad(i\in V(H)).
 \label{eq:paper-positive-dual-gaps}
\end{equation}

For each permitted pattern $P$, choose a witness $j(P)$ with
$P\subseteq N_H^-(j(P))$. Fix a positive integer $K$. The new graph
$D_K$ has two kinds of vertices. For each $j\in V(H)$, it has a
distinguished vertex $c_j$ with label $j$ and pattern $N_H^-(j)$.
For each $P\in\mathcal P(H)$, it also has $K w_P$ vertices with
pattern $P$ and label $j(P)$. Distinct vertices are retained when
patterns or labels coincide. If $\ell(v)$ and $P(v)$ denote the label
and pattern of a vertex, define
\begin{equation}
 v\to u\quad\Longleftrightarrow\quad \ell(v)\in P(u).
 \label{eq:paper-pattern-realization-arcs}
\end{equation}

Every pattern satisfies $P(v)\subseteq N_H^-(\ell(v))$, so
$\ell(v)\notin P(v)$; hence \eqref{eq:paper-pattern-realization-arcs}
creates no loop. If both $v\to u$ and $u\to v$ held, the same
containments would give both label arcs $\ell(v)\to\ell(u)$ and
$\ell(u)\to\ell(v)$ in $H$. For unequal labels this is a digon,
and for equal labels it is a loop. Thus $D_K$ is oriented. Its
distinguished vertices induce exactly $H$.

Consider any source $v$ with label $i$. Its first targets are exactly
the vertices $u$ with $i\in P(u)$. We claim that its exact second
targets are exactly those with $i\in R_H(P(u))$. Indeed, if
$v\to x\to u$ and $i\notin P(u)$, put $j=\ell(x)$. The first
arc and $P(x)\subseteq N_H^-(j)$ imply $i\to j$ in $H$, while
the second arc says $j\in P(u)$. Thus $i\in R_H(P(u))$.
Conversely, if $i\in R_H(P(u))$, choose
$j\in N_H^+(i)\cap P(u)$. The distinguished vertex $c_j$ supplies
the path $v\to c_j\to u$, and $i\notin P(u)$ excludes a first
arc. The target cannot be $v$: since
$P(v)\subseteq N_H^-(i)$, such a $j$ would give a digon $i\leftrightarrow j$
in $H$.

The extra pattern vertices therefore contribute $K\gamma_i$ to the
gap of every source labelled $i$. The distinguished vertices contribute
$g_H(i)$, because they supply exactly the first and exact second
targets of $i$ in $H$. We have proved the identity
\begin{equation}
 g_{D_K}(v)=K\gamma_i+g_H(i)
                    \qquad\text{whenever }\ell(v)=i.
 \label{eq:paper-realization-gap}
\end{equation}
Choose $K$ so large that the right-hand side is positive for every
$i\in V(H)$. This is possible by
\eqref{eq:paper-positive-dual-gaps}. Then $D_K$ is a finite
all-deficient oriented graph. Consequently statement~(1) rules out
infeasibility for every $H$, proving statement~(2).
\end{proof}

\begin{corollary}[An equivalent whole-class formulation]
\label{cor:paper-whole-class-equivalence}
Seymour's second neighborhood conjecture is equivalent to the assertion
that every finite nonempty oriented graph has a maximal incoming class
consisting entirely of Seymour vertices.
\end{corollary}
\begin{proof}
The whole-class assertion plainly implies the conjecture. Conversely,
assume the conjecture, and hence universal pattern feasibility by
Theorem~\ref{thm:pattern-duality}. If each maximal incoming class of
$D$ contained a positive-gap vertex, choose one such representative
from each class. Lemma~\ref{lem:paper-maximal-representatives-covered}
and Theorem~\ref{thm:pattern-transfer} would give a Seymour vertex
among these representatives, a contradiction.
\end{proof}

\subsection{Lifting through iterated representative cores}

The representative operation can be repeated, and it preserves more
than the existence of one Seymour vertex: a certificate at a later
stage lifts to a certificate at every preceding stage.

\begin{lemma}[Lifting a pattern certificate]
\label{lem:core-lifting}
Let $C$ contain one representative from each maximal incoming class of
an oriented graph $H$, and put $K=H[C]$. The extension by zero of any
pattern certificate for $K$ is a pattern certificate for $H$.
\end{lemma}
\begin{proof}
Fix $P\subseteq N_H^-(v)$. Choose a representative $c\in C$ such
that $N_H^-(v)\subseteq N_H^-(c)$. Then
\[
 P\cap C\subseteq N_K^-(c),
 \qquad R_K(P\cap C)\subseteq R_H(P)\cap C.
\]
For the second inclusion, a vertex of $C$ outside $P\cap C$ is
outside $P$, and an arc into $P\cap C$ is also an arc of $H$ into
$P$. If $\lambda$ is the certificate on $C$ and
$\widehat\lambda$ is its zero extension, it follows that
\[
 \widehat\lambda(R_H(P))
 \geq\lambda(R_K(P\cap C))
 \geq\lambda(P\cap C)
 =\widehat\lambda(P).
\]
The lifted vector is nonnegative and nonzero, so it is a certificate.
\end{proof}

\begin{corollary}[Iterated-core sufficient conditions]
\label{cor:paper-iterated-core-sufficient}
Starting from $D$, repeatedly take the induced subgraph on one
representative from each maximal incoming class of the current graph.
If any resulting core has at most five vertices, or is a tournament
of any order, then the original graph $D$ has a Seymour vertex in that
core.
\end{corollary}
\begin{proof}
The small-core certificate lemma or the tournament certificate lemma
supplies a certificate at the specified stage. Apply
Lemma~\ref{lem:core-lifting} successively in reverse order. This gives
a certificate on $D$ supported in the specified core. Applying
Theorem~\ref{thm:pattern-transfer} with the whole graph as its core
gives an original Seymour vertex in that support.
\end{proof}

\begin{corollary}[Reduction to incomparable incoming neighborhoods]
\label{cor:paper-terminal-incoming-core}
Universal pattern feasibility, and therefore Seymour's second neighborhood
conjecture, is equivalent to pattern feasibility for oriented graphs
whose incoming neighborhoods are distinct and pairwise incomparable
under inclusion.
\end{corollary}
\begin{proof}
Whenever the representative operation reduces the vertex set, repeat
it. The process terminates because the order decreases. If the
operation no longer reduces the order, every incoming class is a
singleton and every class is maximal. The incoming neighborhoods are
therefore distinct and pairwise incomparable. A certificate for this
terminal graph lifts to the original graph by
Lemma~\ref{lem:core-lifting}. The reverse implication is immediate.
\end{proof}

This is an induced-subgraph reduction: it deletes vertices and does not
add arcs. Its conclusion concerns the original neighborhoods of vertices
in the final core. For a single chosen chain it does not assert that
an entire original maximal incoming class consists of Seymour vertices;
the whole-class conclusions above use representative selection under
the negation of that stronger assertion.

\section{Predecessor containment and extremal counterexamples}
\label{sec:paper-predecessor-cloning}

In this section we study a surgery that preserves positive gaps by
copying outgoing neighborhoods under incoming-neighborhood containment.
The surgery itself is valid in an arbitrary finite oriented graph.
Its structural consequences will explicitly state whether they require
minimum order or a stronger lexicographic extremal choice.

We call $D$ \emph{all-deficient} if $g_D(v)>0$ for every vertex.
Such a graph is a counterexample to the second neighborhood conjecture.
The gap-one predecessor and gap-one cycle phenomena occur in the
minimal-counterexample analysis of Brantner et al.~\cite{BrantnerEtAl2009}.
Minimum-order reductions also play a central role in
\cite{EspunyDiazEtAl2025}. The arguments below give statements about
entire incoming classes and establish the precise extremal hypotheses
needed for an independent-module normal form.

\subsection{Copying one row and copying all rows}

All containment predicates and all copied neighborhoods in the next
two results are evaluated in the original graph $D$.

\begin{lemma}[One-row copying under predecessor containment]
\label{lem:paper-single-row-copy}
Let $a\ne b$ and suppose that
$N_D^-(a)\subseteq N_D^-(b)$. Replace the outgoing neighborhood of
$a$ by $N_D^+(b)$, leaving every other outgoing neighborhood unchanged,
and denote the resulting graph by $D^{a\leftarrow b}$. Then this graph
is oriented and
\[
 N_{D^{a\leftarrow b}}^+(a)=N_D^+(b),\qquad
 N_{D^{a\leftarrow b}}^{++}(a)=N_D^{++}(b).
\]
For every $v\ne a$,
\[
 N_{D^{a\leftarrow b}}^+(v)=N_D^+(v),\qquad
 N_{D^{a\leftarrow b}}^{++}(v)\subseteq N_D^{++}(v).
\]
In particular, the new gap at $a$ is $g_D(b)$, all other gaps do not
decrease, and an all-deficient graph remains all-deficient.
\end{lemma}
\begin{proof}
There is no original arc $b\to a$: it would put $b$ in $N_D^-(a)$
and hence in $N_D^-(b)$, contrary to the absence of loops. Nor is there
an original two-step path $b\to x\to a$, since $x\to a$ would imply
$x\to b$ and form a digon with $b\to x$. Thus $a$ is neither a first
nor a second target of $b$.

Every $x\in N_D^+(b)$ is distinct from $a$ and cannot point to $a$,
again because $x\to a$ would imply $x\to b$. The new row at $a$
therefore creates neither a loop nor a digon. The intermediate vertices
in $N_D^+(b)$ retain their outgoing neighborhoods. Their union of
outgoing targets contains neither $a$, by the two-step exclusion just
proved, nor $b$, by the absence of digons. Subtracting the copied first
set $N_D^+(b)$ consequently gives exactly $N_D^{++}(b)$ as the new
second neighborhood of $a$.

For $v\ne a$, only a two-step path using the changed row could produce
a new second target. Such a path has the form $v\to a\to x$, where
$v\to a$ and $b\to x$ are original arcs. Incoming containment gives
the original arc $v\to b$, so $v\to b\to x$ was already an original
two-step path. The first set at $v$ is unchanged, so no new exact
second target appears. Some second targets may disappear. The gap
statements now follow by counting the first and second sets.
\end{proof}

\begin{theorem}[A simultaneous predecessor map]
\label{thm:clone-map}
Let $\rho:V(D)\to V(D)$ satisfy
\begin{equation}
 N_D^-(v)\subseteq N_D^-(\rho(v))\qquad(v\in V(D)).
 \label{eq:paper-clone-containment}
\end{equation}
Define $D^\rho$ by simultaneously setting
\[
 N_{D^\rho}^+(v)=N_D^+(\rho(v))\qquad(v\in V(D)).
\]
Then $D^\rho$ is oriented, and for every $v$,
\begin{equation}
 N_{D^\rho}^{++}(v)\subseteq N_D^{++}(\rho(v)),
 \qquad g_{D^\rho}(v)\geq g_D(\rho(v)).
 \label{eq:paper-clone-map-gaps}
\end{equation}
The map $\rho$ need not be injective, and its image may contain vertices
whose rows are also changed. In particular, all-deficiency is preserved.
\end{theorem}
\begin{proof}
A new loop at $v$ would mean an original arc $\rho(v)\to v$.
By \eqref{eq:paper-clone-containment}, this would imply an original
loop at $\rho(v)$. If $v\to u$ and $u\to v$ were both new arcs,
the original arcs $\rho(v)\to u$ and $\rho(u)\to v$, together with
containment at $u$ and $v$, would give the original label arcs
$\rho(v)\to\rho(u)$ and $\rho(u)\to\rho(v)$. These form a digon
when the two images differ and a loop when they coincide. Thus the
new graph is oriented.

Consider a new two-step path $v\to u\to x$. It means that
$\rho(v)\to u$ and $\rho(u)\to x$ are original arcs. Applying
incoming containment at $u$ to the first of these arcs yields the
original arc $\rho(v)\to\rho(u)$. Hence
\[
 \rho(v)\to\rho(u)\to x
\]
is an original two-step path. Its intermediate vertex cannot equal
$\rho(v)$, since that would require a loop, and its endpoint cannot
equal $\rho(v)$, since that would require a digon. If $x$ is an exact
second target of $v$ in $D^\rho$, then
$x\notin N_{D^\rho}^+(v)=N_D^+(\rho(v))$. Therefore it is an
original exact second target of $\rho(v)$, proving the inclusion in
\eqref{eq:paper-clone-map-gaps}. The first neighborhood has exactly
the copied original size, so the gap inequality follows.
\end{proof}

\begin{corollary}[Copying a prescribed subset]
\label{cor:paper-subset-copy}
Let $A\subseteq V(D)$, and choose $\rho(a)$ for each $a\in A$ with
$N_D^-(a)\subseteq N_D^-(\rho(a))$. Extend $\rho$ by the identity
outside $A$. Then simultaneous copying gives
\[
 g_{D^\rho}(a)\geq g_D(\rho(a))\quad(a\in A),\qquad
 g_{D^\rho}(v)\geq g_D(v)\quad(v\notin A).
\]
\end{corollary}
\begin{proof}
The extended map satisfies \eqref{eq:paper-clone-containment} at every
vertex, so Theorem~\ref{thm:clone-map} applies.
\end{proof}

The simultaneous formulation is essential for coupled choices of the
map: all containments refer to one graph, and all rows are copied from
that same graph. A sequence that recomputes containment after each
individual copy is a different procedure and is not needed here.

\subsection{A unit-predecessor cover at minimum order}

For the remainder of this subsection, assume that $D$ is all-deficient
and has the smallest possible number of vertices among all
all-deficient oriented graphs. This is a global minimum-order
assumption. For $f\in V(D)$ define
\[
 A_f=\{a\in V(D):a\to f\text{ and }g_D(a)=1\}.
\]
The next obstruction concerns a simultaneous choice of replacement
vertices outside the set being copied.

\begin{lemma}[Unit predecessors cannot all be copied outside their set]
\label{lem:paper-unit-predecessor-obstruction}
For no vertex $f$ is there a map
\[
 \rho:A_f\longrightarrow V(D)\setminus A_f
 \quad\text{such that}\quad
 N_D^-(a)\subseteq N_D^-(\rho(a))\quad(a\in A_f).
\]
In particular $A_f\ne\varnothing$ for every $f$.
\end{lemma}
\begin{proof}
Suppose such a map exists. Extend it by the identity outside $A_f$
and copy the rows simultaneously. The resulting graph $D^\rho$ is
all-deficient by Theorem~\ref{thm:clone-map}.

We claim that every predecessor of $f$ in $D^\rho$ has gap at least
two. An unchanged predecessor was originally a predecessor outside
$A_f$, and hence had positive integral gap at least two; its gap has
not decreased. If a changed vertex $a\in A_f$ still points to $f$,
then its chosen image $b=\rho(a)$ originally pointed to $f$.
Since $b\notin A_f$, we again have $g_D(b)\geq2$, and the copy
inequality gives $g_{D^\rho}(a)\geq2$. There are no other new
predecessors because all rows outside $A_f$ were unchanged.

Delete $f$ from $D^\rho$. Each of its predecessors loses one first
target and can only lose exact second targets among the remaining
vertices; its gap therefore remains at least one. Every other
remaining vertex loses no first target and acquires no exact second
target, so its gap also remains positive. The deleted graph is
nonempty: an all-deficient oriented graph has positive outdegree at
every vertex and therefore at least three vertices. We have produced
a smaller all-deficient graph, contradicting minimum order.

If $A_f$ were empty, the empty map would meet the stated conditions
and the same deletion argument would apply.
\end{proof}

\begin{theorem}[Every vertex receives a maximal pure-unit class]
\label{thm:unit-class-cover}
For every $f\in V(D)$, some maximal incoming class $C$ satisfies
\[
 C\subseteq A_f.
\]
Equivalently, every member of $C$ has original gap one and sends an
original arc to $f$.
\end{theorem}
\begin{proof}
Suppose no maximal incoming class is contained in $A_f$. Each maximal
class then has a representative outside $A_f$. For every $a\in A_f$,
choose a maximal class whose incoming set contains $N_D^-(a)$, and
let $\rho(a)$ be its chosen representative outside $A_f$. Such a
maximal class exists in the finite inclusion poset. The resulting map
satisfies
\[
 N_D^-(a)\subseteq N_D^-(\rho(a)),\qquad
 \rho(a)\notin A_f\qquad(a\in A_f),
\]
contradicting Lemma~\ref{lem:paper-unit-predecessor-obstruction}.
\end{proof}

Call a maximal incoming class \emph{pure-unit} if all its members have
gap one, and let $\mathcal C_1$ be the family of these classes.

\begin{corollary}[A complete-arc cycle of pure-unit classes]
\label{cor:paper-unit-class-cycle}
There are at least three classes in $\mathcal C_1$. Define an oriented
quotient on $\mathcal C_1$ by placing an arc $C\to C'$ when every
vertex of $C$ points to every vertex of $C'$. This quotient has positive
minimum indegree and contains a directed cycle of length at least three.
Selecting one vertex from each class on such a cycle gives an original
directed cycle all of whose vertices have gap one.
\end{corollary}
\begin{proof}
Theorem~\ref{thm:unit-class-cover} makes $\mathcal C_1$ nonempty.
Fix $C'\in\mathcal C_1$ and $f\in C'$. The theorem gives
$C\in\mathcal C_1$ with every member of $C$ pointing to $f$.
Every vertex of $C'$ has the same incoming neighborhood as $f$;
hence all these arcs extend to complete arcs from $C$ to $C'$.
The classes are distinct, since an incoming class is independent.
Thus each quotient vertex has an incoming arc.

The quotient has no loops or digons, by its definition and the
orientedness of $D$. Following predecessors in this finite quotient
produces a directed cycle. Such a cycle has at least three vertices.
Its complete interclass arcs allow any choice of one original vertex
from each cycle class, proving the last assertion.
\end{proof}

Combining this with Corollary~\ref{cor:paper-counterexample-six-incoming-classes},
a minimum-order counterexample would have at least six maximal incoming
classes, at least three of which are pure-unit. Neither assertion says
that every maximal incoming class must be pure-unit.

\subsection{A lexicographic extremal normal form}

Suppose a counterexample exists. Among all all-deficient oriented graphs,
choose $D$ according to the following successive objectives:
\begin{enumerate}
 \item minimize $|V(D)|$;
 \item subject to the first objective, minimize $|A(D)|$;
 \item subject to the first two objectives, maximize
       $\sum_{v\in V(D)}g_D(v)$;
 \item subject to the first three objectives, maximize
       \[
        \Psi(D)=\sum_F
              |\{v\in V(D):N_D^+(v)=F\}|^2,
       \]
       where the sum ranges over the distinct outgoing neighborhoods.
\end{enumerate}
These choices exist conditionally on the existence of a counterexample:
the first minimum is a positive integer, and after fixing it only finitely
many labelled oriented graphs remain. The second objective ranges over
all counterexamples of that order, rather than just subgraphs of one
chosen graph. Call a graph selected in this way \emph{globally extremal}.

\begin{lemma}
\label{lem:paper-extremal-connectivity}
A globally extremal counterexample is strongly connected and minimal
under deletion of any nonempty set of arcs while the vertex set is fixed.
\end{lemma}
\begin{proof}
If it were not strongly connected, choose a sink strongly connected
component. It is a proper nonempty vertex subset with no outgoing arc
to its complement. Both first and second neighborhoods of its vertices
therefore remain unchanged in the induced subgraph. That subgraph
would be all-deficient, contrary to minimum order. If deleting a
nonempty arc set preserved all-deficiency, the resulting graph would
have the same order and fewer arcs, contrary to the second objective.
\end{proof}

\begin{proposition}[Degree and gap monotonicity under incoming containment]
\label{prop:normal-degree}
In a globally extremal counterexample, if $a\ne b$ and
$N_D^-(a)\subseteq N_D^-(b)$, then
\[
 d_D^+(a)\leq d_D^+(b).
\]
If these outdegrees are equal, then $g_D(a)\geq g_D(b)$. In particular,
equal incoming neighborhoods imply equal outdegrees and equal gaps.
If in addition $a\to b$, the missing degrees satisfy
\[
 \mu_D(a)\geq\mu_D(b)+1,
 \qquad \mu_D(v)=|V(D)|-1-d_D^+(v)-d_D^-(v).
\]
\end{proposition}
\begin{proof}
By Lemma~\ref{lem:paper-single-row-copy}, copying the outgoing row
of $b$ to $a$ preserves all-deficiency. It changes the arc count by
$d_D^+(b)-d_D^+(a)$. A negative change would contradict the second
extremal objective, which proves the degree inequality.

If the degrees coincide, the copied graph has the same order and arc
count. Its new gap at $a$ is $g_D(b)$, and all other gaps do not
decrease. Its total gap is thus at least the old total plus
$g_D(b)-g_D(a)$. The third objective excludes a positive increase,
giving $g_D(a)\geq g_D(b)$. When the incoming neighborhoods are
equal, apply both conclusions in both directions.

If $a\to b$, then $a\in N_D^-(b)\setminus N_D^-(a)$, so incoming
containment yields $d_D^-(b)\geq d_D^-(a)+1$. Therefore
\[
 \mu_D(a)-\mu_D(b)
 =d_D^+(b)-d_D^+(a)+d_D^-(b)-d_D^-(a)\geq1.
\]
\end{proof}

\begin{theorem}[Independent incoming classes in an extremal graph]
\label{thm:twin-normal-form}
In a globally extremal counterexample,
\[
 N_D^-(a)=N_D^-(b)\quad\Longrightarrow\quad N_D^+(a)=N_D^+(b).
\]
Consequently every incoming class is an independent module: each outside
vertex has uniform adjacency and orientation to the entire class.
Members of a class have the same first and exact second neighborhoods,
the same gap, and the same missing degree.
\end{theorem}
\begin{proof}
Suppose that $a$ and $b$ have the same incoming neighborhood but
different outgoing neighborhoods $F_a$ and $F_b$.
Proposition~\ref{prop:normal-degree} gives equal outdegrees and equal
gaps. Both single-row copying directions are permitted. They preserve
the first two extremal objectives; their total gaps cannot decrease,
and by the third objective cannot increase. Thus both operations also
preserve the third objective.

Let $r_a$ and $r_b$ be the numbers of vertices whose outgoing
neighborhoods are $F_a$ and $F_b$, respectively. Interchange the names
if necessary so that $r_a\leq r_b$, and copy $b$'s row to $a$.
Every other outgoing row is unchanged. The value of the fourth
objective changes by
\[
 (r_a-1)^2+(r_b+1)^2-r_a^2-r_b^2
       =2(r_b-r_a)+2>0.
\]
This remains the correct formula if the old row class disappears.
The increase contradicts the fourth extremal objective, proving
equality of the outgoing neighborhoods.

An incoming class is independent, as noted above. Equality of both
incoming and outgoing neighborhoods gives uniform adjacency to every
outside vertex, including uniform nonadjacency when no arc is present.
The unions of outgoing neighborhoods over the common first set are
identical for all members of the class. Such a union contains no
member of the class: a path $a\to x\to b$ between two class members
would imply $x\to a$ from their equal incoming sets and hence a
digon. Subtracting the common first set therefore gives identical
exact second sets. Equality of the gaps and missing degrees follows.
\end{proof}

The independent-module conclusion concerns the selected globally
extremal graph. The row-copying inequalities hold without that selection,
but the degree ordering, gap ordering, and final symmetrization use its
successive objectives. In particular, arc-deletion minimality alone
does not supply the tie-breaking properties used in
Proposition~\ref{prop:normal-degree} and
Theorem~\ref{thm:twin-normal-form}.

\section{Protection closure and median feeds}\label{sec:paper-closure}
The incoming-core reduction takes induced subgraphs. We now keep the vertex
set fixed and add arcs. Every convenience predicate in this section refers
to the graph present at the stated stage.

\subsection{Two safe operations}
A missing direction $a\to b$ is \emph{forced} in $H$ if some vertex $w$
satisfies $a\triangleleft_H w$ and $w\to b\in A(H)$.

\begin{lemma}\label{lem:paper-forced-step}
Forced directions are consistent, uniquely convenient in $H$, and mandatory
in every tournament completion preserving all protections of $H$.
Adding all forced directions simultaneously preserves every old protection.
\end{lemma}
\begin{proof}
For a witness $a\triangleleft_H w$, $w\to b$, every predecessor $p$ of $a$
also points to $w$, and hence reaches $b$ within two steps. Thus $a\to b$
is convenient. Its reverse fails at $w$, since $w$ cannot reach its protected
far head $a$. Opposite directions cannot both be forced: each would have
been proved convenient while the other witness proves it nonconvenient.
In a tournament, $a\triangleleft w$ also gives outgoing domination, so
$w\to b$ forces $a\to b$.

Fix an old $y\triangleleft_H x$. An old predecessor of $y$ already points
to $x$. For a newly incoming $a\to y$, choose its witness
$a\triangleleft_H w$, $w\to y$. Then $w\to x$ originally. If $ax$ is
missing it is forced $a\to x$ in the same round. If $x\to a$ were an old
arc, protection at $a,w$ would give $x\to w$, a digon with $w\to x$.
Thus the required new incoming nesting is preserved in every case.
\end{proof}

There is a second, sequential operation. If $ab$ is missing,
$N_H^-(a)\subseteq N_H^-(b)$, and $b$ has no protection successor in $H$,
add the single arc $a\to b$. Call this a \emph{weak step}.

\begin{lemma}\label{lem:paper-weak-step}
A weak step is convenient in its preceding graph, preserves all old
protections, and creates $a\triangleleft b$.
\end{lemma}
\begin{proof}
Every old predecessor of $a$ already points to $b$. Only $N^-(b)$ changes.
An old relation $y\triangleleft b$ becomes easier to satisfy, and an old
relation $b\triangleleft x$ is excluded by the hypothesis on $b$.
All other incoming comparisons are unchanged. The added arc together with
the old containment gives $a\triangleleft b$.
\end{proof}

Starting at $D$, repeatedly perform a full forced round; when no forced
direction remains, perform one available weak step and recompute.
Stop when neither operation is available, obtaining a partial graph $Q$.
Any deterministic choice among eligible weak steps may be used.
The process terminates after at most $M(D)$ new pairs. Protections acquired
at any stage persist. In particular, every added arc has a tail with a
protection successor in $Q$. Simultaneous arbitrary weak choices are not
part of the construction.

\subsection{A fixed partial graph and a complete reference family}
Since $Q$ has no remaining forced pair, each $y\triangleleft_Q x$ has
outgoing as well as incoming nesting. Indeed, $x\to z$ excludes $z\to y$
by incoming nesting; if $yz$ were missing, $y\to z$ would still be forced.
A two-position swap therefore shows that every global median order of $Q$
extends protection: reversing $x,y$ gains their mutual arc and loses no
contribution from an intermediate vertex. Missing arcs contribute zero.

Complete every remaining missing pair forward in a linear extension of
the protection order of $Q$. Every such tournament preserves protection.
To see this, consider $y\triangleleft_Q x$ and a newly incoming $a\to y$.
The reference has $a<y<x$. An old $x\to a$ would have forced $y\to a$,
contrary to terminality and the missing pair $ay$. Thus $a\to x$ is either
old or is also completed forward. This proves the incoming nesting.

Let $T_*$ maximize the ordinary median score over the family of
\emph{all} these extensions. A remaining pair of $Q$ is called free.

\begin{theorem}[Rigidity and original neighbourhoods]\label{thm:paper-feed-invariance}
Every global median of $T_*$ respects protection in $Q$ and puts every free
arc forward. At any feed $f$ of such an order,
\[
 N_{T_*}^+(f)=N_Q^+(f)=N_D^+(f),\qquad
 N_Q^{++}(f)=N_D^{++}(f).
\]
If $r$ is the number of free pairs, the maximum tournament median score is
$\med(Q)+r$, and the possible feeds across the maximizing completions are
exactly the median feeds of $Q$.
\end{theorem}
\begin{proof}
Tournament domination makes every median respect protection. If a free arc
were backward in one median of $T_*$, use that median as a new reference.
This fixes the arcs of $Q$ and turns all backward free arcs forward,
strictly increasing the score. This contradicts maximality and applies to
every median of the same maximizing tournament. The argument requires the
entire extension family.

A feed has no outgoing free arc, since it would be backward. Nor can it be
the tail of an arc added during closure: its persistent protection successor
would have to occur later. Thus its first set was unchanged at every stage.
Induct over the stages. Any new two-step path at $f$ uses an original first
arc and a newly added second arc. Convenience in the preceding graph says
that this original predecessor $f$ already reached the head within two steps.
A first neighbour stays first throughout, so no new exact second head occurs.
Old exact second paths persist because the first set is fixed.

A median of $Q$ is an admissible reference and scores $\med(Q)+r$ after
completion. No order in any reference tournament can exceed this number:
its $Q$ contribution is at most $\med(Q)$ and its free contribution at most
$r$. Equality forces both bounds to be attained, proving the score and feed
assertions.
\end{proof}

\subsection{Equality sedimentation and the original far budget}
For a median order $L$ of a tournament with feed $f$, a vertex
$y\in N^-(f)$ is good if some $a\in N^+(f)$ occurs before $y$ in $L$ and
has $a\to y$. Let $G_L$ be the good set. The strong feed theorem of
Havet and Thomass\'e gives
\[
 |N^+(f)|\le |G_L|\le |N^{++}(f)|.
\]
When $|G_L|=|N^+(f)|$, their global sedimentation lemma replaces $L$ by
the order consisting of its bad vertices, then $f$, then $N^+(f)\cup G_L$,
preserving each block's previous relative order. It remains a global median
of the same tournament~\cite{HavetThomasse2000}. Each vertex of the last
block moves strictly later. These statements concern ordinary global medians.

\begin{theorem}[Strict original surplus]\label{thm:paper-strict-surplus}
Let $f$ be a median feed of $T_*$ with $g_D(f)>0$, and set
\[
 A_f=N_{T_*}^{++}(f)\setminus N_D^{++}(f),\qquad
 R_f=\mathcal M_Q(f),\quad m_f=|R_f|.
\]
Then $m_f>0$ and $t_D(f)\ge g_D(f)+1$.
If $g_{T_*}(f)=0$, then
\[
 \varnothing\ne R_f\subseteq U_{T_*}(f)\cap
       \bigl(\mathcal F_D(f)\setminus A_f\bigr),\qquad
 t_D(f)\ge g_D(f)+m_f.
\]
No all-deficiency or minimality assumption on the other vertices is needed.
\end{theorem}
\begin{proof}
Original protection survives. By Theorem~\ref{thm:paper-feed-invariance},
\[
 N_{T_*}^{++}(f)=N_D^{++}(f)\,\dot\cup\,A_f,
 \quad A_f\subseteq\mathcal F_D(f),\quad
 g_{T_*}(f)=g_D(f)-|A_f|\le0.
\]
If $R_f$ were empty, $f$ would be whole in $Q$. Every far head of a whole
vertex is protected in that graph. Preservation of these protections would
prevent any further second head at $f$, giving $g_{T_*}(f)=g_D(f)>0$.
This contradiction proves $m_f>0$.

If the completed gap is zero, strong feed equality gives
$G_L=N_{T_*}^{++}(f)$. A partner $r\in R_f\cap G_L$ has a free incoming
arc $r\to f$, since the first set at $f$ is original. Sedimentation puts
$f$ before $r$ while keeping a median, contradicting free-arc rigidity.
Thus every remaining partner is far in $T_*$. It was originally missing,
so it cannot be an originally protected head, and old exact second paths
survive. Hence it lies in $\mathcal F_D(f)\setminus A_f$. The two disjoint
sets have sizes $m_f$ and $g_D(f)$. If the completed gap is negative instead,
integrality gives $|A_f|\ge g_D(f)+1$ directly.
\end{proof}

\begin{theorem}[An original partner cycle]\label{thm:paper-original-cycle}
Under the zero-gap hypotheses of Theorem~\ref{thm:paper-strict-surplus},
every vertex of $Q[R_f]$ has an incoming neighbour there, and the original
graph $D[R_f]$ contains a directed cycle. Consequently
\[
 m_f\ge3,\qquad t_D(f)\ge g_D(f)+m_f\ge g_D(f)+3.
\]
\end{theorem}
\begin{proof}
The feed is maximal in terminal protection. For $r\in R_f$, terminality
of the weak operation gives $p\in N_Q^-(r)\setminus N_Q^-(f)$.
Since $r$ is far from $f$ in $T_*$, $p\to r$ implies $p\to f$ there.
Neither direction on $pf$ was present in $Q$, so $p\in R_f$. Thus each
remaining partner has a predecessor in the induced terminal graph.

Fix the final set $R_f$ throughout the construction. For a forced internal
addition $a\to b$, let $a\triangleleft_H w$, $w\to b$ be its preceding
witness. Far-ness of $b$ in $T_*$ gives $w\to f$ there. If $wf$ were adjacent
in $Q$, its direction would already be $w\to f$, and the retained protection
$a\triangleleft w$ would force the still-missing $a\to f$. Therefore
$w\in R_f$. The internal addition shortcuts an old internal two-step path.
All simultaneous such shortcuts preserve acyclicity of an acyclic induced
graph.

For a single weak internal addition $a\to b$, creating the first cycle
would require an old internal path $b\leadsto p\to a$. It cannot have length
one, because $ab$ was missing. Old incoming containment gives $p\to b$,
already completing an old cycle. Thus weak steps cannot create the first
cycle either. An acyclic $D[R_f]$ would remain acyclic, whereas a finite
graph of positive minimum indegree has a cycle. Orientedness excludes cycle
lengths one and two.
\end{proof}

The stronger $+3$ bound applies to zero completed gap only. Omitting all weak
steps still gives Theorem~\ref{thm:paper-strict-surplus}, but does not give
the partner-cycle conclusion. The original cycle in the theorem is not a
claim that every terminal cycle lifts edge by edge, nor that its vertices
are actual bad-direction sources.

\begin{corollary}\label{cor:paper-low-far}
Every oriented graph with $t_D(v)\le1$ at every vertex has a Seymour vertex.
More generally this holds if $t_D(v)\le g_D(v)$ at every deficient vertex.
\end{corollary}
\begin{proof}
Otherwise every original gap is positive and any feed contradicts the
strict-surplus theorem.
\end{proof}

If $h_D(f)$ counts the protected original far heads, partitioning the other
vertices gives $t_D(f)-g_D(f)=n-1-2d_D^+(f)-h_D(f)$. Hence every
original-positive-gap closure feed satisfies
\begin{equation}\label{eq:paper-degree-localization}
 2d_D^+(f)+h_D(f)\le n-2.
\end{equation}
At $n=2\delta^+(D)+3$, such a feed has minimum original outdegree and at
most one protected far head. This locates a possibly deficient feed; it
does not assert that a minimum-degree vertex already has the SNP in $D$.

\subsection{The negative-gap equality branch}
Every new original-far head belongs to $G_L$: its second arc cannot lie in
$Q$, and therefore is free and forward, putting its original first intermediate
before the head. Write
$b_L=|N_D^{++}(f)\setminus G_L|$ and
$s_L=|G_L|-d_D^+(f)$. Then
\begin{equation}\label{eq:paper-good-slack}
 s_L=|A_f|-g_D(f)-b_L\ge0.
\end{equation}
When $s_L=0$, sedimentation excludes $R_f\cap G_L$. Its second partners
number at most $|A_f|-g_D(f)$ and are original nongood second vertices;
its far partners number at most $t_D(f)-|A_f|$. Thus
\begin{equation}\label{eq:paper-equality-partners}
 |R_f|\le t_D(f)-g_D(f)\qquad(s_L=0).
\end{equation}
This includes negative completed gap. Put
$R_2=R_f\cap N_D^{++}(f)$, $R_\infty=R_f\cap U_{T_*}(f)$, and
$U=R_\infty\setminus\Reach_{D[R_f]}(R_2)$, where reachability includes
the seed vertices. The supplement proves that closure creates no new
reachability from $R_2$ into $R_\infty$ inside $R_f$. If the whole set $U$
is nonempty, its original induced graph $D[U]$ contains a directed cycle
(Theorem~\ref{S-thm:weak-equality-original-partner-path-cycle}).
When $s_L>0$, the cited equality-sedimentation move is unavailable. This
does not assert that every median order of the tournament has the same slack.

\section{Residual consequences in minimal counterexamples}\label{sec:paper-residuals}
Throughout this section, unless stated otherwise, $D$ is in the MIN setting
of Section~\ref{sec:paper-notation}. Put
\[
 \eta(v)=\mu(v)-1+d^{--}(v)-d^{++}(v),\qquad
 R=\sum_v\eta(v)=2M-n.
\]
The last identity counts exact directed distance-two pairs by their two
endpoints. Unlike the core certificate, the inequalities below use original
arc-set minimality.

\subsection{The original boundary and partner estimates}
We collect the local estimates used here. Full proofs are given in the
supplement so that the external-boundary convention is explicit.

\begin{lemma}[Original deletion estimates]\label{lem:paper-packet}
Every original gap belongs to $\{1,2\}$. If $\partial S\ne\varnothing$,
then $|\partial^2S|<|\partial S|$. Moreover $\eta(v)\ge0$ for every $v$,
and
\[
 \eta(v)=2\mu(v)+g(v)-1-|P(v)|.
\]
Fix $v$ with $P(v)\ne\varnothing$, and write
\[
 \begin{gathered}
 C=\partial P(v),\quad X=\partial^2P(v),\quad J=N^{--}(v)\setminus C,\\
 B=J\cap\mathcal M(v),\quad k=|B|,\quad s=|J\cap N^+(v)|.
 \end{gathered}
\]
Delete the $s$ arcs from $v$ into $J$, obtaining $\widehat D$. Then
\[
 N_{\widehat D}^{++}(v)\subseteq
 E=(\mathcal M(v)\setminus B)\,\dot\cup\,X,
 \qquad L=E\setminus N_{\widehat D}^{++}(v)
\]
satisfies
\begin{equation}\label{eq:paper-packet-bound}
 |L|\le\eta(v)-2k-g(v)\ind_{\{s>0\}},\qquad
 \bigcup_{r\in\mathcal M(v)\setminus B}(N^+(r)\cap U(v))\subseteq L.
\end{equation}
Every exceptional partner $r\in B$ admits the convenient direction
$r\to v$. Residual zero gives $t(v)=0$, residual one gives $t(v)\le1$,
and $\eta(v)=t(v)=1$ gives $P(v)\ne\varnothing$ and
$N^{--}(v)=\partial P(v)$.
\end{lemma}
\begin{proof}
The gap bound follows by deleting one outgoing arc: other first sets are
unchanged and their second sets can only shrink, while the deleted tail's
gap decreases by at most two. The external-boundary deletion argument is
Lemma~\ref{S-lem:boundary}; the original target partition yields
Theorem~\ref{S-thm:main} and the residual identity. The shared remainder
estimate is Lemma~\ref{S-lem:partner-head-budget}; it subtracts the original
second degree if $s=0$, and the second degree forced by minimality after the
nonempty deletion if $s>0$. Its union bound counts distinct heads.
Zero-residual protection and the unit bound are proved in the supplementary
local-budget section. At residual one, a nonempty deletion gives $t=0$;
otherwise equality forces $J$ empty. The empty-$P$ identity gives a whole
vertex and $t=0$, proving the final assertion when $t=1$.
\end{proof}

The set $L$ concerns an original deletion. It is never recomputed in a
closure graph and then assigned the original minimality hypothesis.
In particular, an exceptional missing far head can lie in $X$; no inclusion
$X\subseteq N^-(v)$ is assumed.

\subsection{A residual-four threshold for zero completed gap}
Let $S_v=\mathcal M(v)\cap U(v)$. For nonempty $P(v)$,
\begin{equation}\label{eq:paper-missing-far}
 |S_v|\le k+|L|\le\eta(v)-k-g(v)\ind_{\{s>0\}}.
\end{equation}
Indeed, a nonexceptional missing far head belongs to $E$ and cannot become
reachable when arcs are deleted, so it lies in $L$. When $k=0$, in fact
$\mathcal F(v)\subseteq L$. A far incoming head $y$ that is not protected
has a predecessor $r$ not pointing to $v$. It cannot be an outneighbour of
$v$, since that would reach $y$ in two steps. Thus it is a missing partner,
and the union bound in \eqref{eq:paper-packet-bound} applies.

\begin{theorem}[Zero-gap residual threshold]\label{thm:paper-eta-four}
At a zero-gap feed $f$ of the weak-closure completion, let
$m=|\mathcal M_Q(f)|$. Then
\[
 \eta_D(f)\ge m+1\ge4.
\]
If $\eta_D(f)=4$, its original packet necessarily satisfies
\[
 P_D(f)\ne\varnothing,\quad m=3,\quad k=s=0,\quad g_D(f)=1,
 \quad t_D(f)=4.
\]
Furthermore $|P_D(f)|=2\mu_D(f)-4\ge2$,
$N_D^{--}(f)=\partial P_D(f)$, $L=\mathcal F_D(f)$,
and $|\partial^2P_D(f)|=|\partial P_D(f)|-1$.
\end{theorem}
\begin{proof}
Theorem~\ref{thm:paper-original-cycle} gives $m\ge3$,
$\mathcal M_Q(f)\subseteq S_f$, and $t_D(f)\ge g_D(f)+m$.
If $P(f)$ is empty, the residual identity gives
$\eta(f)=2\mu(f)+g(f)-1\ge2m+g(f)-1\ge m+1$.
For nonempty $P(f)$ and $k>0$, \eqref{eq:paper-missing-far} gives
$\eta(f)\ge m+k+g(f)\ind_{\{s>0\}}\ge m+1$.
For $k=0$, the inclusion $\mathcal F(f)\subseteq L$ gives
$\eta(f)\ge g(f)+m+g(f)\ind_{\{s>0\}}\ge m+1$.

At residual four these bounds force $P(f)\ne\varnothing$, $m=3$,
$s=0$ and $k\le1$. If $k=1$, all three original missing far heads are
exactly the terminal partners, the unique exceptional partner is among
them, and the other two already consume the bound $|L|\le2$.
Their original induced graph has a directed cycle, hence a triangle.
The exceptional vertex receives an original arc from a nonexceptional
partner on this triangle. The shared head bound puts it in $L$ too,
a contradiction. Thus $k=0$; now $g+m\le t\le4$ gives $g=1,t=4$.
The residual identity gives the stated size of $P$. The partition gives
$N^{--}=\partial P$, and $\mathcal F\subseteq L$, $|L|\le4$ gives
$L=\mathcal F$. Here $\widehat D=D$ and
$|E|=d^{++}(f)+4=\mu(f)+|\partial^2P|$; combine this with
$\eta=\mu-1+|\partial P|-d^{++}=4$ to obtain the final equality.
\end{proof}

This threshold does not apply to strictly negative completed gaps. At strong
Good equality the supplementary original path theorem gives a finer statement:
if $\eta(f)=2$, every terminal partner is an original nongood second vertex;
if $\eta(f)=3$, there is at most one terminal far partner and it receives an
original arc from such a second partner. The unique far partner at residual
three may be exceptional. These conclusions leave both the all-second case
and the strictly positive Good-slack case unresolved.

\subsection{Actual failure sources and the residual margin}
For a missing direction $a\to b$, its original failure set is
$\{x:x\to a,\ b\in U_D(x)\}$. Let $\mathcal A$ be the vertices occurring
in the failure set of a direction on a bad missing pair. Opposite directional
witnesses are nonadjacent. The cross-witness graph joins every such pair of
witnesses. A missing pair is flexible if it is bad or convenient in both
directions.

\begin{theorem}[One uncontrolled actual source]\label{thm:paper-one-exception}
In an arbitrary oriented graph, suppose some vertex $u$ has the following
properties: every actual source in $\mathcal A\setminus\{u\}$ has $t\le1$,
and every cross-witness pair between two such sources is flexible. Then the
graph has a Seymour vertex. In MIN, at least two distinct actual failure
sources have residual at least two. In particular every positive residual
allocation $(m,1^k)$ is excluded for $m\ge2$ and finite $k$.
\end{theorem}
\begin{proof}
The complete argument is Theorem~\ref{S-thm:one-exception-short-general}.
Its first optimization uses all canonical protection-reference completions,
maximizing the forward score and then the position of $u$. A unit-budget
feed is balanced. An opposite witness other than $u$ makes a flexible
incoming arc backward after sedimentation, whereas witness $u$ moves $u$
strictly later. Thus the selected feed must be $u$, which therefore has
no original protection successor. A separate fixed completion avoiding
bad failures at $u$, followed by a median maximizing its position, yields
the second contradiction. The fixed-cover argument is proved independently
in Theorem~\ref{S-thm:protected-independent-heavy-cover}.

In MIN, an actual source has positive residual; residual one has $t\le1$,
and the unit--unit opposite-witness pair is bad. If at most one actual
source had residual at least two, the general result would apply. The
allocation consequence follows even when inactive residuals are unrestricted.
\end{proof}

\begin{theorem}[Parity-sensitive missing-pair margin]\label{thm:paper-residual-seven}
Every MIN graph satisfies $R=2M-n\ge7$, with $R\ge8$ at even order.
Every arc-set-minimal all-deficient graph, whether strong or not, satisfies
\[
 M\ge\left\lceil n/2\right\rceil+3+\ind_{\{n\text{ even}\}}.
\]
If $o$ vertices lie outside a sink strong component, the right side can be
increased by $o$. Equality in the uniform bound implies strong connectivity.
\end{theorem}
\begin{proof}
The supplementary two-support theorem gives at least three positive
residual vertices (Corollary~\ref{S-cor:three-positive-residual-support}).
Theorem~\ref{thm:paper-one-exception} requires two residuals at least two.
Hence, at total residual at most six, the only possible positive partitions are
\[
 (2,2,1),\quad(3,2,1),\quad(2,2,2),\quad(2,2,1,1).
\]
The first two are excluded by the general $(m,2,1)$ theorem
(Theorem~\ref{S-thm:no-three-two-one}); the last two by
Theorems~\ref{S-thm:eta222-excluded} and \ref{S-thm:no-two-two-one-one}.
These are complete graph arguments, not an enumeration of possible orders.
Thus $R\ge7$; parity gives $R\ge8$ for even $n$.

For a sink component of order $k$ with $o=n-k$, the inherited minimality
and the outside missing-pair count in Theorem~\ref{S-thm:sink} give
\[
 M(D)\ge\left\lceil k/2\right\rceil+3+\ind_{\{k\text{ even}\}}
               +\left\lceil o(k-2)/2\right\rceil.
\]
An all-deficient graph has minimum outdegree at least two, so $k\ge5$.
Using $k-2\ge3$ and checking the two parities yields at least
$\lceil n/2\rceil+o+3+\ind_{\{n\text{ even}\}}$.
If $o>0$ this is strictly stronger than the uniform bound.
\end{proof}

Thus at $n=2\delta+3$ one has $M\ge\delta+5$, and at $n=2\delta+4$
one has $M\ge\delta+6$. The remaining residual-seven partitions after
these filters are
\[
 (3,3,1),\ (3,2,2),\ (3,2,1,1),\ (2,2,2,1),\ (2,2,1,1,1).
\]
Their weak-closure median feeds must have strictly negative completed gap,
by Theorem~\ref{thm:paper-eta-four}. None of these five partitions is
excluded by that observation alone.

\subsection{An additional global constraint on a near-whole feed}
Suppose a weak-closure feed $f$ has one original missing partner $r$.
Terminal whole-feed exclusion leaves $r$ still missing. Absence of a weak
$r\to f$ step supplies a current predecessor $p$ of $r$ not pointing to
$f$. Since $f$ is adjacent to every other vertex, $f\to p$; exact feed
neighbourhood invariance gives $r\in N_D^{++}(f)$.
If also $P_D(f)$ is empty, each unprotected original far head $y$ must
receive an original arc from $r$, the only possible missing-partner witness
against protection. Put $O=N_D^+(f)$, $d=d_D^+(f)$, $g=g_D(f)$ and $t=t_D(f)$.
Then $\mathcal F_D(f)\subseteq\partial_D^2O$. Every member of $O$ reaches
$f$ originally within two steps, so $f\in\partial_D^2O$ as well. Therefore
\[
 t+1\le|\partial_D^2O|<|\partial_D O|=d-g,
 \qquad d\ge g+t+2,\qquad n\ge g+3t+5.
\]
This uses the global original boundary inequality. It is stronger than
the constraints obtained by deleting only arcs with tail $f$, but still
leaves unbounded parameter ranges.

\section{Exact certificates and reproducibility}\label{sec:paper-verification}
The computational component proves a finite statement about cores. It does
not enumerate all original graphs covered by Theorem~\ref{thm:core-five};
their total order is unbounded. The mathematical transfer theorem passes
from the finite certificates to those graphs.

\subsection{Enumeration and the verification boundary}
For a labelled core on $k$ vertices, unordered pairs are listed in
lexicographic order. A ternary digit encodes absence, the lower-to-higher
direction, or its reverse, with the first pair at the least significant digit.
This enumerates all $3^{\binom{k}{2}}$ labelled oriented cores.
For each core $H$, the certificate is a nonzero nonnegative integer vector $w$.
For every set $P$ with a common outneighbour, the verifier checks
\[
 \sum_{i\notin P:\ N_H^+(i)\cap P\ne\varnothing}w_i
       \ge\sum_{i\in P}w_i.
\]
Dividing by $\sum_iw_i$ gives the normalized vector used in the proof.

\begin{table}[htbp]
\centering
\caption{Complete incoming-pattern certificate verification.}
\label{tab:paper-certificate-counts}
\begin{tabular}{@{}rrrr@{}}
\toprule
Core order & Labelled cores & Nonempty inequalities & Largest weight\\
\midrule
1 & 1 & 0 & 1\\
2 & 3 & 2 & 1\\
3 & 27 & 54 & 1\\
4 & 729 & 3\,078 & 1\\
5 & 59\,049 & 459\,270 & 3\\
\midrule
Total & 59\,809 & 462\,404 & 3\\
\bottomrule
\end{tabular}
\end{table}

Including one empty pattern for each core gives $522\,213$ checks.
At order five, $59\,009$ certificates are uniform on their support;
the other forty use weights from $\{0,1,3\}$. All forty corresponding
cores are tournaments, also covered by the analytic density argument.
An exact rational solver discovered these forty vectors, but all stored
vectors are checked directly with integer arithmetic. Solver correctness
is not part of the verifier's trusted boundary.

The independently organized replays reconstruct adjacency and external
predecessor sets from the ternary code. One uses Boolean matrices rather
than the generator's incoming-bitmask unions. They check coverage, vector
length, nonnegativity, nonzero mass, and every permitted pattern. Integer
sums in the supplied tables are at most nine; the principal replay also
checks that all sums are safe integers.

\begin{lstlisting}[caption={Independent certificate verification, pseudocode.},label={alg:paper-verifier}]
for k = 1,...,5:
    assert table[k].length == 3^(k*(k-1)/2)
    for code = 0,...,table[k].length-1:
        H = decode_oriented_graph(k, code)
        w = table[k][code]
        assert w is a nonzero vector of nonnegative integers
        for P subset of V(H):
            if some j receives an arc from every vertex of P:
                R = {i outside P: i points to some vertex of P}
                assert sum(w[i] for i in R) >= sum(w[i] for i in P)
\end{lstlisting}

The supplied verification directory can be checked with standard Node.js,
without a solver, network access or third-party packages:
\begin{lstlisting}
node verify_small_core_certificates.js
node verify_core_pattern_certificates_independent.js
\end{lstlisting}
The SHA-256 checksum of the order-five table is
\begin{center}\small\ttfamily
c9fe139b76a3866a7bf17e408d12f344823fa\linebreak
1a54fca306fd656139e601bfd5a
\end{center}
The line break is typographical. Reproducing the generator is unnecessary
to validate a supplied certificate.

\subsection{Controls for the completion arguments}
Finite graph checks supplement the closure proofs but do not establish
their universal quantifiers. The following controls test their scope.
In an adjacency mask, bit $j$ of the $i$th integer means $i\to j$.

The six-vertex graph $(4,4,24,32,32,3)$ has weak closure
$(6,4,24,48,32,3)$ and a maximizing completion
$(30,28,56,48,32,3)$. Its unique median is $(0,1,2,3,4,5)$.
At feed $5$, the original gap is one, the completed gap is minus one,
$t_D=2$, and the terminal missing degree is one. Thus the general $+1$
surplus is attained in the negative-gap branch. Good slack is one, so
equality sedimentation cannot move the order.

The eleven-vertex original masks
\[
 (6,8,8,17,1793,80,144,48,33,65,129)
\]
give terminal partners $\{5,6,7\}$ at feed $0$, with original cycle
$5\to6\to7\to5$. Exact subset dynamic programming verifies partial
median score $23$ and maximizing tournament score $49$ for the reference
$(8,9,10,5,6,7,1,2,3,4,0)$. Here
$g_D(0)=1$, $g_T(0)=0$ and $t_D(0)=4$, attaining the zero-gap $+3$ bound.
Both controls contain Seymour vertices and are not counterexamples to
the conjecture or to the MIN residual estimates.

A ten-vertex control also satisfies all outgoing-star deletion conditions
at one feed, with local residual two, yet has stable negative completed
gap. Its original first boundary and external second boundary both have
size three, violating the global boundary lemma. This distinguishes
single-feed criticality from full original arc-set minimality. Its explicit
data and the exact-score programs are supplied with the supplementary records.

\subsection{Other computations and their scope}
Copying and closure programs test identities on arbitrary graphs, including
several changed rows and shared heads. They do not simulate a hypothetical
all-deficient minimal graph. Searches that fail to find a graph or a difficult
branch are not proofs of nonexistence. Earlier minimum-degree-seven branch
logs are archived as local computational records and are not used to infer
a complete minimum-degree theorem. The exact five-core theorem has the
exhaustive finite coverage specified above, independently of those searches.

\section{Further consequences and unresolved steps}\label{sec:paper-discussion}
The incoming-pattern formulation isolates a precise global question: does
every finite oriented core admit a pattern certificate? Theorem~\ref{thm:pattern-duality}
shows that universal feasibility would settle Seymour's conjecture, while
an infeasible core would yield a concrete counterexample. The present tables
establish only the core-order-five range. Iterated incoming-core reduction
removes repeated or contained incoming sets from the remaining certificate
problem, but does not prove feasibility on all larger terminal cores.

The copying and completion methods impose different constraints. In a
minimum-order counterexample the maximal gap-one classes cover every target;
a further lexicographic choice makes incoming classes true independent
modules. A closure feed can nevertheless remain originally deficient because
it gains several second neighbours. The zero-gap branch has an original
partner cycle and residual at least four. The negative-gap branch can have
strict Good slack or an original nongood second partner. We have not proved
that the original unit-class cover eliminates all these possibilities.

Two supplementary results delimit other approaches. Cyclic three-copy
amplification sends order $n$ and minimum outdegree $\delta$ to
\[
 n'=3n,\qquad\delta'=n+\delta,\qquad n'-2\delta'=n-2\delta,
\]
while preserving vertex gaps, missing degrees and local residuals.
A fixed near-half minimum-degree theorem for all oriented graphs would
therefore imply the conjecture. Amplification does not preserve arc-set
minimality, so a minimal residual bound cannot simply be transferred to
the amplified dense graph.

For a fixed protection-preserving tournament completion $T$, set
$\Delta=g_T-g_D$, and let $N_T$ be the signed second-minus-first matrix.
The exact correction budget has dual form
\[
 \max_{\substack{\lambda\ge0,\ \mathbf1^T\lambda=1\\N_T^T\lambda\ge0}}
       \lambda^T\Delta
   =\min_{w\ge0}\max_x\bigl(\Delta_x+(N_Tw)_x\bigr).
\]
An explicit six-class clone family has optimum $-k/3$ for every canonical
reference completion, even after strong dense amplification. These graphs
already have Seymour vertices. They refute an unconditional completion-budget
claim, not the conjecture. The supplement supplies a common integer dual
certificate for the whole family. A normalized correction strictly greater
than $-1$ would suffice for an all-deficient graph; its universal existence
remains unproved.

The specialized $n=2\delta+3$ results in the supplement include exact
deletion identities, minimum-degree predecessor estimates, capacity and
heavy-set constraints, and conditional degree thresholds. Their capacity
variable $\sigma$ is distinct from $\eta$. Earlier finite pure-unit bounds
and one-heavy exclusions are recorded as consequences of the unbounded
source theorem rather than as independent contributions.

\paragraph{Verification and attribution.}
The arguments and programs underwent independent internal reconstruction
and replay. This is not a proof-assistant certificate or an external referee
report. AI tools assisted research, drafting and implementation; human
authorship, review and approval are required before submission. Literature
comparison establishes the cited inputs and their scope, but does not certify
priority for every auxiliary observation. The source, certificate tables,
dependency inventory and verification records support further examination.
No complete proof of the general conjecture or actual counterexample is asserted.

\bibliographystyle{plain}
\bibliography{../references}
\end{document}
